\section{本章小结}
\label{sec:ls-normed-inner-product-summary}

本章回答了章首的三个问题：两个向量相差多少，一列近似是否趋于极限，两个方向是否垂直。前两个问题由范数回答，第三个问题需要内积。

范数给向量规定大小，从而可以度量误差。常用坐标范数之间满足$\lVert\bm x\rVert_q\leq\lVert\bm x\rVert_p\leq d^{1/p-1/q}\lVert\bm x\rVert_q$（$p\leq q$），比较常数依赖维数，所以同一个误差在不同范数下的读数不同。

有了范数就能谈收敛与连续。向量列与矩阵列的收敛等价于逐坐标、逐元素收敛；映射的联合极限要求所有趋近方式给出同一结果，只检查若干路径只能否定、不能证明极限存在。闭有界集上的连续函数取到最值并且一致连续，闭性、有界性与连续性三个条件缺一不可。

有限维结构带来更强的结论。任意两个范数等价，所以收敛、连续与柯西性质都不依赖所选范数，变的只是误差的数值和比较常数。有限维赋范空间完备，是Banach空间，子空间都是闭集；绝对收敛的向量与矩阵级数收敛，并有尾和估计。诱导矩阵范数给出线性映射的最大放大倍数，有限维线性映射因而都Lipschitz连续；对偶范数把输入误差的限制转化为线性读取的最大变化；商范数衡量忽略一个子空间中的差别以后剩余的大小。一般范数下向子空间的最近点存在，但可能不唯一。

内积进一步给出长度、夹角和正交。内积诱导的长度是范数，但并非每个范数都来自内积，平行四边形恒等式可以检验这一点；有限维内积空间是Hilbert空间。Gram矩阵记录一组向量的两两内积，它总是半正定的，并且恰在这组向量线性无关时正定；反过来，每个半正定矩阵都是某组向量的Gram矩阵，这又推出逐元素乘积保持半正定。

正交分解把向量拆成子空间中的部分与垂直的部分，正交投影是内积诱导距离下唯一的最近点；矩阵是正交投影，当且仅当它幂等且对称。标准正交基使坐标可以由内积直接读出，并保持长度；Gram--Schmidt正交化保证这样的基总存在，QR分解把这一过程写成矩阵形式。线性泛函的内积表示与伴随，把上一章的代数对偶与本章的几何结构联系起来：在标准正交基下伴随就是转置，在一般基下还要计入两端的Gram矩阵。

这些结论都依赖有限维：范数等价、完备性与最近点存在的证明，都用到了有限个坐标或闭有界集上的最值。它们为后续各章提供共同依据。第~\ref{chap:eigen-spectral-decomposition}章在正交换基下寻找线性作用的简单表示，正交矩阵与自伴算子在那里尤为重要；矩阵幂和级数使用收敛与次乘性，迭代使用完备性和误差估计，多元微分把局部变化表示为线性映射并用算子范数控制余项。
